% docs/manual/user/chapters/09-migration.tex
% 第 9 章：从旧版迁移

\chapter{从旧版迁移}

\openbench{} v3.0 把 v2.0 的多文件 JSON/NML 配置（main + ref + sim + variable defs，共 4-6 个文件）整合为单一 \file{openbench.yaml}。本章介绍迁移工具 \cli{openbench migrate} 的用法，以及如何核对迁移结果。

\section{v2.0 vs v3.0 配置差异}

\subsection{文件数}

\begin{itemize}
  \item v1.x：Fortran NML 风格，主 namelist + 多个子 namelist（5+ 文件）
  \item v2.0：JSON / 早期 YAML，main 配置 + ref 配置 + sim 配置 + 变量定义（4-6 文件）
  \item v3.0：单一 \file{openbench.yaml}（可选用 \texttt{!include} 拆分，但不强制）
\end{itemize}

\subsection{字段命名}

v3.0 在某些字段上做了重命名：

\begin{longtable}{l l l}
  \toprule
  v2.0 字段 & v3.0 字段 & 说明 \\
  \midrule
  \endhead
  \texttt{syear} / \texttt{eyear} & \texttt{years: [\textit{s},\textit{e}]} & 年份合并为列表 \\
  \texttt{options.data\_root} & \texttt{reference.data\_root} & 移到 reference 块 \\
  \texttt{options.\textit{X}} & \texttt{project.\textit{X}} & 旧 options 块全部并入 project \\
  \texttt{compare\_tim\_res} & \texttt{project.tim\_res} & 全局目标分辨率 \\
  \texttt{compare\_grid\_res} & \texttt{project.grid\_res} & 全局目标分辨率 \\
  \texttt{IGBP\_groupby} & \texttt{project.IGBP\_groupby} & 同名只是位置改了 \\
  \bottomrule
\end{longtable}

旧字段在 \openbench{} v3.0 中通过 loader 兼容层自动迁移到新位置；运行时会打印 deprecation 警告。

\section{迁移命令：\cli{openbench migrate}}

\subsection{基本用法}

\begin{minted}{bash}
# 从 v2.0 JSON 主配置
openbench migrate old_config.json
# 默认输出 openbench.yaml

# 指定输出文件名
openbench migrate old_config.json -o my_evaluation.yaml

# 从 v1.x Fortran NML
openbench migrate main.nml -o openbench.yaml
\end{minted}

\subsection{支持的输入格式}

\begin{itemize}
  \item \texttt{.json}：v2.0 主配置
  \item \texttt{.yaml} / \texttt{.yml}：早期 YAML 配置（v2.0 后期）
  \item \texttt{.nml}：v1.x Fortran namelist
\end{itemize}

\subsection{输出示例}

\begin{exampleBox}[迁移成功]
\begin{minted}{text}
✓ Read 5 config files
✓ 12 evaluation variables
✓ 3 simulation models
✓ Written to openbench.yaml

Next: openbench check openbench.yaml
\end{minted}
\end{exampleBox}

迁移工具会：

\begin{enumerate}
  \item 读主配置（JSON / NML / 早期 YAML）
  \item 跟踪并解析 referenced 文件（ref namelist、sim namelist、variable defs）
  \item 把所有信息合并为单一 \file{openbench.yaml}
  \item 应用字段重命名规则
  \item 检测并填入合理默认（如 \yamlkey{tim_res}、\yamlkey{grid_res}）
\end{enumerate}

\section{典型迁移流程}

\subsection{步骤 1：备份旧配置}

\begin{minted}{bash}
cp -r old_config_dir old_config_dir.backup
\end{minted}

\subsection{步骤 2：跑 migrate}

\begin{minted}{bash}
openbench migrate old_config_dir/main.nml -o openbench.yaml
\end{minted}

\subsection{步骤 3：审阅生成的 yaml}

\begin{minted}{bash}
less openbench.yaml
\end{minted}

重点核对：

\begin{itemize}
  \item \yamlkey{project.years} 是不是预期的 \texttt{[start, end]}
  \item \yamlkey{evaluation.variables} 数量与列表
  \item \yamlkey{reference} 每个变量都映射到了 v3.0 注册表里的有效 reference 名（旧 v2.0 的"自由文本"reference 名可能在 v3.0 已被规范化）
  \item \yamlkey{simulation} 各 entry 的 \yamlkey{root_dir} 路径在新机器上仍正确
\end{itemize}

\subsection{步骤 4：跑 check}

\begin{minted}{bash}
openbench check openbench.yaml
\end{minted}

如果 check 报错，常见原因：

\begin{itemize}
  \item Reference 名在新 catalog 中变了（如 v2.0 \texttt{GLEAM4} → v3.0 \texttt{GLEAM\_v4.2a\_LowRes}）→ 手动改 \yamlkey{reference} 那一行
  \item 旧路径在新环境不存在 → 改 \yamlkey{simulation.<label>.root_dir}
  \item 旧字段 v3.0 已废弃 → 删掉
\end{itemize}

\subsection{步骤 5：跑 dry-run 对比}

\begin{minted}{bash}
openbench run openbench.yaml --dry-run
\end{minted}

\openbench{} 会列出将要评估的所有 (sim, ref, var) 三元组。与你印象中的旧配置任务清单核对。

\subsection{步骤 6：实际运行并对比结果}

如果有旧配置的评估结果（v2.0 跑过的），用 v3.0 重跑同一份迁移后配置，理论上结果应当数值一致：

\begin{itemize}
  \item Metrics 数值应当与旧版一致到 6 位小数（浮点边角差异可接受）
  \item Score 一致
  \item 图像视觉一致
\end{itemize}

\warninline{\textbf{已知差异}：v3.0 的 metric 计算路径走新 \modname{core/metrics.py}（与 v2.0 完全相同的算法），但若你是从 v1.x Fortran 迁移过来，部分 metric 的 NaN 处理细节可能略有差异（v3.0 在 NaN 处理上更严格）。如果数值差异大于浮点边角，请先用相同变量的少量数据 cross-check，再下结论。}

\section{无法迁移的字段}

某些 v2.0 字段在 v3.0 已被弃用或重设计：

\subsection{自定义 Fortran 模块路径}

v1.x / v2.0 早期允许在配置中引用自定义 Fortran 处理脚本。v3.0 全部用 Python，不读这些路径。

\subsection{硬编码 reference 路径}

v2.0 在配置里硬编码 reference 数据集路径（每条 entry 写完整路径）。v3.0 把 reference 路径搬到 catalog（\file{reference\_catalog.yaml}），配置只引用名字。迁移工具会保留旧路径作为 \yamlkey{simulation.<label>.root_dir} 的注释，但 reference 路径需要你确认 catalog 里已注册。

\subsection{某些 v1.x 输出布局选项}

v3.0 的输出目录结构是固定的（详见第 7 章）。v1.x 可配的 \texttt{output\_subdir\_layout} 等选项被忽略。

\section{保留旧配置作为参考}

迁移工具不会删除旧配置文件。建议：

\begin{itemize}
  \item 旧配置目录改名为 \file{old\_config\_v2/}（或类似）保留
  \item 新 \file{openbench.yaml} 与旧配置共存几次评估周期，便于回退
  \item 长期项目档案中保留旧配置便于历史追溯
\end{itemize}

\section{何时迁移可以跳过}

如果你是 \openbench{} 新用户，没有 v2.0 的旧配置，直接用 \cli{openbench init} 生成全新配置即可，本章可忽略。

如果旧配置非常简单（只评估 1-2 个变量），手动重写比迁移工具更快、更干净。

\section{下一步}

下一章是 FAQ 与排错指南：安装、配置、数据、运行、远程各方面常见问题的诊断流程。
